\documentclass[aspectratio=169]{beamer}
\usetheme{Madrid}
\usecolortheme{whale}

\usepackage[utf8]{inputenc}
\usepackage[spanish]{babel}
\usepackage{amsmath,amssymb}
\usepackage{graphicx}
\usepackage{listings}
\usepackage{booktabs}

\lstset{
    language=Python,
    basicstyle=\ttfamily\footnotesize,
    keywordstyle=\color{blue}\bfseries,
    stringstyle=\color{red},
    commentstyle=\color{gray}\itshape,
    breaklines=true,
    showstringspaces=false
}

\title[INT210 - Sesión 0.4]{Sesión 0.4: Bucles Indeterminados (\texttt{while}) y Convergencia}
\subtitle{Módulo 0: Fundamentos Algorítmicos en Python}
\author[Diplomado en Ciencia de Datos]{Cátedra de Inteligencia Artificial y Ciencia de Datos}
\institute[INT210]{Machine Learning From Scratch}
\date{Semana 1}

\begin{document}

\frame{\titlepage}

\begin{frame}{Contenido de la Sesión}
    \tableofcontents
\end{frame}

\section{Introducción y Analogía}

\begin{frame}{La Analogía Infantil: La Tetera y el Bucle Indeterminado}
    \begin{columns}
        \column{0.55\textwidth}
        \begin{itemize}
            \item \textbf{Condición Dinámica:} Mantener fuego encendido mientras la tetera NO silbe.
            \item \textbf{Señal Centinela:} El silbato dispara el apagado inmediato (\texttt{break}).
            \item \textbf{Límite de Seguridad:} Si pasa 1 hora sin silbido, apagar por prevención de incendios (\textit{Guard Timeout}).
        \end{itemize}
        \column{0.45\textwidth}
        \begin{alertblock}{Iteración Indeterminada}
            Se emplea cuando el número total de ciclos depende de una condición emergente en tiempo de ejecución.
        \end{alertblock}
    \end{columns}
\end{frame}

\section{Convergencia Numérica y la Parada de Turing}

\begin{frame}{El Problema de la Parada y Criterio de Convergencia $\epsilon$}
    \begin{columns}
        \column{0.5\textwidth}
        \textbf{El Teorema de Parada (Turing, 1936):}
        \begin{itemize}
            \item Imposibilidad matemática de predecir si un bucle general terminará.
            \item Mandato de ingeniería: implementar siempre \texttt{max\_iter}.
        \end{itemize}
        \column{0.5\textwidth}
        \textbf{Criterio de Parada Numérica:}
        \begin{equation*}
            \Delta_k = |x_{k+1} - x_k| < \epsilon
        \end{equation*}
        \begin{itemize}
            \item Base fundamental del Machine Learning (Descenso de Gradiente).
            \item Tolerancia típica: $\epsilon = 10^{-4}$.
        \end{itemize}
    \end{columns}
\end{frame}

\section{Implementación en Ciberseguridad y Arte}

\begin{frame}[fragile]{Network Listener Continuo con Centinela}
\begin{lstlisting}
while socket_activo:
    paquete = socket.recibir()
    if paquete.flag == 'FIN_ACK':
        cerrar_sesion_tcp()
        break
    procesar_payload(paquete)
\end{lstlisting}
    \vspace{0.2cm}
    \textbf{Renderizado Adaptativo de Iluminación:}
\begin{lstlisting}
while delta_energia > epsilon and rebotes < MAX_REBOTES:
    energia_acumulada += rebote_actual
    rebote_actual *= reflectancia
    rebotes += 1
\end{lstlisting}
\end{frame}

\begin{frame}{Resumen de la Semana 1}
    \begin{center}
    \begin{tabular}{lll}
        \toprule
        \textbf{Sesión} & \textbf{Estructura Clave} & \textbf{Dominio Aplicado} \\
        \midrule
        \textbf{0.1} & \texttt{if-elif-else} & Firewall de red y Luminancia RGB \\
        \textbf{0.2} & \texttt{match-case} / Dict Dispatch & Triaje SOC $\mathcal{O}(1)$ y Render multiformato \\
        \textbf{0.3} & \texttt{for} / Comprensiones & Matrices de imagen y Auditoría de puertos \\
        \textbf{0.4} & \texttt{while} / Centinelas & Network Listeners y Convergencia $\epsilon$ \\
        \bottomrule
    \end{tabular}
    \end{center}
    \vspace{0.3cm}
    \textbf{¡Felicitaciones! Has completado la nivelación algorítmica de la Semana 1.}
\end{frame}

\end{document}
